\documentclass[11pt]{article}
\usepackage[margin=0.72in]{geometry}
\usepackage[T1]{fontenc}
\usepackage[utf8]{inputenc}
\usepackage[spanish,es-tabla]{babel}
\usepackage{lmodern}
\usepackage{microtype}
\usepackage{amsmath,amssymb,amsfonts}
\usepackage{booktabs}
\usepackage{array}
\usepackage{longtable}
\usepackage{enumitem}
\usepackage{listings}
\usepackage{xcolor}
\usepackage[hidelinks]{hyperref}

\definecolor{chlblue}{HTML}{1F77B4}
\definecolor{chlgreen}{HTML}{2CA02C}
\definecolor{chlamber}{HTML}{FFB000}
\definecolor{chlred}{HTML}{D62728}
\definecolor{chldark}{HTML}{202332}
\definecolor{chlmuted}{HTML}{5E6678}
\definecolor{codebg}{HTML}{F6F8FA}

\hypersetup{
  pdftitle={CHL: memoria de agente local-first, MCP-native y aprendible},
  pdfauthor={Documento técnico de desarrollo},
  pdfsubject={Memoria semántica ligera, MCP, LSH, Hypervector, HDC, VSA},
  pdfkeywords={CHL, MCP, memoria semántica, SimHash, LSH, HDC, VSA, benchmark}
}

\setlength{\parindent}{0pt}
\setlength{\parskip}{0.52\baselineskip}
\emergencystretch=2em
\setlist[itemize]{leftmargin=1.35em,itemsep=0.15em,topsep=0.2em}
\setlist[enumerate]{leftmargin=1.55em,itemsep=0.15em,topsep=0.2em}
\renewcommand{\arraystretch}{1.15}

\lstdefinestyle{chlcode}{
  backgroundcolor=\color{codebg},
  basicstyle=\ttfamily\footnotesize,
  keywordstyle=\color{chlblue}\bfseries,
  commentstyle=\color{chlmuted},
  stringstyle=\color{chlgreen!70!black},
  numbers=left,
  numberstyle=\tiny\color{chlmuted},
  stepnumber=1,
  numbersep=8pt,
  frame=single,
  framerule=0.3pt,
  rulecolor=\color{chlmuted!35},
  breaklines=true,
  tabsize=2,
  showstringspaces=false,
  captionpos=b
}
\lstset{style=chlcode}

\newcommand{\CHL}{\textsc{CHL}}
\newcommand{\LSH}{\textsc{LSH}}
\newcommand{\HDC}{\textsc{HDC}}
\newcommand{\VSA}{\textsc{VSA}}
\newcommand{\MCP}{\textsc{MCP}}
\newcommand{\simh}{\operatorname{sim}_{\mathrm H}}
\newcommand{\hd}{d_{\mathrm H}}
\newcommand{\popcnt}{\operatorname{popcnt}}
\newcommand{\xor}{\oplus}
\newcommand{\bits}{\{0,1\}}
\newcommand{\pmone}{\{-1,+1\}}

\title{\vspace{-1.3em}\textbf{\CHL{}: memoria de agente local-first, \MCP{}-native y aprendible}\\
\Large Arquitectura final de memoria semántica ligera para IA}
\author{Documento técnico de desarrollo\\Adaptado al estado final del sistema}
\date{Mayo de 2026}

\begin{document}
\maketitle

\begin{abstract}
Este documento describe la versión final de \CHL{} como una memoria de agente local-first, \MCP{}-native, persistente y aprendible. La idea central sigue siendo usar firmas binarias sensibles a la localidad para recuperar vecinos semánticos, pero el sistema final ya no depende de un ``hash'' criptográfico como representación de significado. En su lugar, combina: normalización y canonicalización, un lexicón aprendido de conceptos y frases, firmas semánticas, hipervectores, persistencia en disco, journal incremental, backup binario comprimido y exposición por HTTP y \MCP{}.

La implementación real se reparte entre \texttt{src/memory.js}, \texttt{src/native.js}, \texttt{native/chl\_addon.cc}, \texttt{src/concepts.js}, \texttt{src/server.js} y \texttt{src/mcp.js}. El servidor HTTP expone inspección y exportación del estado y del lexicón; el servidor \MCP{} expone herramientas y recursos para que un LLM lea, consulte y actúe sobre la memoria. Los benchmarks reproducibles en \texttt{scripts/evaluate.js} muestran dos regímenes de interés: en \texttt{bench:large} con 12\,000 documentos y 6\,000 consultas, \texttt{chl-native} alcanza recall@1 de 0.9392 y p95 de 1255.92\,\(\mu s\); en \texttt{bench:huge} con 50\,000 documentos y 20\,000 consultas, el recall@1 global queda en 0.80925 y la latencia p95 en 863.54\,\(\mu s\), con una degradación esperable de paráfrasis al escalar.
\\
El resultado posiciona a \CHL{} no como una base vectorial genérica, sino como una memoria conversacional viva para agentes: inspectable, exportable, persistente y con vocabulario aprendido entre sesiones.
\end{abstract}

\textbf{Palabras clave:} memoria de agente, \MCP{}, \LSH{}, SimHash, \HDC{}, \VSA{}, lexicón aprendido, distancia de Hamming, memoria semántica ligera.

\tableofcontents
\newpage

\section{Introducción}

El objetivo de \CHL{} no es sustituir a los Transformers ni competir como base de datos vectorial universal. El objetivo es más concreto: ofrecer a un agente una memoria local que pueda recordar, inspeccionar, exportar y reutilizar conocimiento con baja latencia y sin infraestructura pesada.

La motivación técnica es doble. Por un lado, los hashes criptográficos destruyen similitud: son excelentes para integridad, pero malos para semántica. Por otro lado, las memorias densas o vectoriales generalistas suelen ser potentes, pero no siempre son la solución más simple cuando lo que se necesita es memoria viva de agente, explicable y conectada al modelo por \MCP{}.

\CHL{} resuelve esa tensión con una arquitectura híbrida:
\begin{itemize}
  \item firmas semánticas compactas para direccionamiento rápido;
  \item hipervectores para composición y re-ranking;
  \item lexicón aprendido para paráfrasis y frases recurrentes;
  \item persistencia local, journal y backup binario;
  \item API HTTP y \MCP{} para inspección y control.
\end{itemize}

\section{Estado final del sistema}

La implementación final ya existe en el repositorio y se divide en módulos concretos:

\begin{table}[h]
\centering
\small
\begin{tabular}{p{0.25\linewidth}p{0.66\linewidth}}
\toprule
\textbf{Archivo} & \textbf{Responsabilidad real} \\
\midrule
\texttt{src/memory.js} & Motor JS de memoria asociativa, representaciones múltiples, scoring y fallback funcional. \\
\texttt{native/chl\_addon.cc} & Addon nativo C++ con hash semántico, búsqueda rápida, persistencia y métricas. \\
\texttt{src/native.js} & Wrapper de alto nivel, fallback JS, persistencia de journal y sidecars del lexicón. \\
\texttt{src/concepts.js} & Canonicalización, lexicón de conceptos y frases, aprendizaje y serialización TSV. \\
\texttt{src/server.js} & API HTTP para estado, backup, restore, lexicón e inspección. \\
\texttt{src/mcp.js} & Herramientas y recursos \MCP{} para memoria, backup, restore, lexicón y estado. \\
\texttt{scripts/evaluate.js} & Benchmarks reproducibles: small, large y huge. \\
\texttt{test/chl.test.js} & Tests de recall, persistencia, HTTP, \MCP{} y lexicón. \\
\bottomrule
\end{tabular}
\caption{Mapa funcional del sistema final.}
\end{table}

\subsection{Capas funcionales}

\begin{enumerate}
  \item \textbf{Memoria local:} entradas, snapshot, journal, buckets, backup binario.
  \item \textbf{Memoria aprendida:} conceptos y frases persistentes en \texttt{memory.concepts.tsv} y \texttt{memory.phrases.tsv}.
  \item \textbf{Interfaz de agente:} herramientas y recursos \MCP{} para acceso directo desde el LLM.
  \item \textbf{Interfaz de inspección:} endpoints HTTP para consultar estado, entradas, backup y lexicón.
\end{enumerate}

\subsection{Perfil real de ejecución}

El runtime soporta dos presets de producto:
\begin{itemize}
  \item \texttt{default}: 128 bits, 32 band bits, 256 dimensiones hipervectoriales, 4096 entradas máximas.
  \item \texttt{large}: 128 bits, 32 band bits, 256 dimensiones hipervectoriales, 20\,000 entradas máximas y 96 candidatos máximos.
\end{itemize}

El benchmark \texttt{bench:huge} no introduce un perfil nuevo de runtime: usa el preset \texttt{large} como estrés sobre un corpus de 50\,000 documentos.

\section{Fundamentos técnicos}

\subsection{Por qué no usar un hash criptográfico como semántica}

Un hash criptográfico ofrece una propiedad opuesta a la que necesitamos: entradas parecidas producen salidas no correlacionadas. Eso es correcto para integridad, pero incorrecto para memoria semántica. En \CHL{}, el hash útil debe preservar vecindad de forma probabilística.

\subsection{Locality-Sensitive Hashing}

Una familia \(\mathcal{H}\) es sensible a localidad si entradas cercanas colisionan con mayor probabilidad que entradas lejanas:
\[
\Pr_{h \sim \mathcal{H}}[h(x)=h(y)] \ge p_1 \quad \text{si } d(x,y)\le r,
\]
\[
\Pr_{h \sim \mathcal{H}}[h(x)=h(y)] \le p_2 \quad \text{si } d(x,y)\ge cr,
\]
con \(p_1 > p_2\) y \(c>1\). Esta es la base conceptual de la recuperación por buckets en \CHL{}.

\subsection{SimHash}

En SimHash, cada bit se define por el signo de una proyección:
\[
h_i(x)=\mathbb{1}[\langle r_i,\phi(x)\rangle \ge 0].
\]
La distancia de Hamming se calcula con XOR y popcount:
\[
\hd(a,b)=\popcnt(a \xor b).
\]
La similitud binaria puede expresarse como:
\[
\simh(a,b)=1-\frac{\hd(a,b)}{m}.
\]

\subsection{HDC y VSA}

\HDC{} y \VSA{} aportan las operaciones de composición:
\begin{itemize}
  \item \textbf{binding}: asociación de roles y valores;
  \item \textbf{bundling}: superposición de múltiples señales;
  \item \textbf{permutación}: codificación de orden y contexto temporal;
  \item \textbf{cleanup}: limpieza por prototipo o memoria asociativa.
\end{itemize}

La implementación final no promete que un hash aislado ``contenga conocimiento''; el conocimiento aparece al combinar firma, hipervector, payload, lexicón aprendido y memoria recuperable.

\section{Arquitectura de memoria}

\subsection{Qué guarda cada entrada}

Cada entrada almacenada por el sistema puede contener:
\begin{itemize}
  \item texto original;
  \item representaciones normalizadas y canonicalizadas;
  \item firma semántica compacta;
  \item hipervector;
  \item payload recuperable;
  \item metadata como calidad, timestamps y contador de acceso;
  \item relaciones o información estructurada cuando el payload lo permita.
\end{itemize}

\subsection{Scoring final}

El ranking recupera candidatos a partir de la firma semántica y los reordena con señales de similitud, hipervector, concepto, recencia, calidad y negación. La forma conceptual del score es:
\[
score(q,e)=\lambda_h \, sim_h + \lambda_v \, sim_v + \lambda_c \, sim_c + \lambda_q \, quality + \lambda_r \, recency + \lambda_n \, negation.
\]

La implementación está repartida entre el fallback JS y el addon nativo para mantener comportamiento alineado.

\subsection{Lexicón aprendido}

El sistema final añadió una capa lateral de aprendizaje: un lexicón persistente de conceptos y frases. Esta capa vive en \texttt{src/concepts.js}, se escribe como TSV y se carga desde las variables de entorno:
\begin{itemize}
  \item \texttt{CHL\_CONCEPTS\_PATH}
  \item \texttt{CHL\_PHRASES\_PATH}
\end{itemize}

En la práctica, esto permite que un LLM o una integración externa reutilice el vocabulario aprendido entre sesiones sin reconstruirlo desde cero.

\subsection{Persistencia}

La persistencia se resuelve en tres niveles:
\begin{enumerate}
  \item \textbf{journal incremental}: eventos de \texttt{remember} y \texttt{learn};
  \item \textbf{backup binario}: archivo comprimido con cabecera \texttt{CHLB} y compresión \texttt{deflate};
  \item \textbf{sidecars del lexicón}: TSV para conceptos y frases.
\end{enumerate}

\section{Interfaz HTTP}

El servidor HTTP expone operaciones de inspección, exportación y restauración. Los endpoints reales son:

\begin{longtable}{p{0.28\linewidth}p{0.62\linewidth}}
\toprule
\textbf{Endpoint} & \textbf{Función} \\
\midrule
\texttt{GET /health} & Estado del servidor con snapshot. \\
\texttt{GET /snapshot} & Snapshot actual de la memoria. \\
\texttt{GET /profile} & Perfil activo (\texttt{default} o \texttt{large}). \\
\texttt{GET /state} & Estado completo, incluyendo entradas, bucket stats y journal. \\
\texttt{GET /entries} & Lista completa de entradas. \\
\texttt{GET /journal} & Journal de mutaciones. \\
\texttt{GET /backup} & Backup JSON compatible. \\
\texttt{GET /backup.bin} & Backup binario comprimido. \\
\texttt{GET /lexicon} & Lexicón aprendido en JSON con conteos y rutas. \\
\texttt{GET /lexicon.concepts.tsv} & Exportación manual de conceptos. \\
\texttt{GET /lexicon.phrases.tsv} & Exportación manual de frases. \\
\texttt{GET /lexicon.tsv} & Exportación combinada para inspección manual. \\
\texttt{POST /remember} & Inserta memoria. \\
\texttt{POST /recall} & Consulta vecinos. \\
\texttt{POST /infer} & Inferencia de mejor respuesta. \\
\texttt{POST /learn} & Aplica feedback. \\
\texttt{POST /restore} & Restaura backup JSON. \\
\texttt{POST /restore.bin} & Restaura backup binario. \\
\bottomrule
\caption{API HTTP real del sistema.}
\end{longtable}

\section{Interfaz \MCP{}}

El diseño \MCP{} es central en el resultado final porque el agente puede acceder a la memoria como herramientas y recursos estándar. Esto encaja con la idea oficial de \MCP{}: exponer recursos para aportar contexto y herramientas para ejecutar acciones controladas \cite{mcpResources,mcpTools}.

\subsection{Herramientas expuestas}

\begin{longtable}{p{0.28\linewidth}p{0.62\linewidth}}
\toprule
\textbf{Tool} & \textbf{Función} \\
\midrule
\texttt{chl\_remember} & Guardar memoria. \\
\texttt{chl\_recall} & Recuperar candidatos. \\
\texttt{chl\_infer} & Inferir la mejor respuesta. \\
\texttt{chl\_learn} & Aplicar feedback. \\
\texttt{chl\_backup} & Exportar backup JSON. \\
\texttt{chl\_backup\_binary} & Exportar backup binario en base64. \\
\texttt{chl\_restore} & Restaurar backup JSON. \\
\texttt{chl\_restore\_binary} & Restaurar backup binario. \\
\texttt{chl\_snapshot} & Inspeccionar snapshot. \\
\texttt{chl\_profile} & Consultar perfil activo. \\
\texttt{chl\_state} & Ver estado completo. \\
\texttt{chl\_entries} & Ver entradas completas. \\
\texttt{chl\_journal} & Ver journal. \\
\texttt{chl\_bucket\_stats} & Ver estadísticas de buckets. \\
\texttt{chl\_clear} & Vaciar memoria. \\
\texttt{chl\_lexicon} & Inspeccionar lexicón aprendido. \\
\texttt{chl\_lexicon\_export} & Exportar lexicón en formato TSV para reutilización manual. \\
\bottomrule
\caption{Herramientas \MCP{} reales de la implementación.}
\end{longtable}

\subsection{Recursos expuestos}

\begin{longtable}{p{0.28\linewidth}p{0.62\linewidth}}
\toprule
\textbf{Resource URI} & \textbf{Contenido} \\
\midrule
\texttt{chl://memory} & Estado completo de memoria. \\
\texttt{chl://profile} & Perfil activo. \\
\texttt{chl://state} & Estado, entradas, bucket stats y journal. \\
\texttt{chl://entries} & Entradas completas. \\
\texttt{chl://journal} & Journal. \\
\texttt{chl://backup} & Backup JSON. \\
\texttt{chl://backup.bin} & Backup binario codificado en base64. \\
\texttt{chl://lexicon} & Lexicón aprendido en JSON. \\
\texttt{chl://lexicon.concepts} & TSV de conceptos. \\
\texttt{chl://lexicon.phrases} & TSV de frases. \\
\texttt{chl://lexicon.tsv} & Exportación TSV combinada. \\
\bottomrule
\caption{Recursos \MCP{} reales de la implementación.}
\end{longtable}

\section{Backup y restauración}

El backup binario usa una estructura compacta con cabecera y compresión \texttt{deflate}. La restauración reconstruye la memoria y, si el backup contiene lexicón, vuelve a escribir los sidecars correspondientes. Esto permite reinicios limpios y reutilización entre sesiones.

La serialización de lexicón también es manualmente exportable en TSV para inspección, sincronización o despliegue en otra instancia.

\section{Implementación nativa y fallback}

El addon nativo (\texttt{native/chl\_addon.cc}) implementa el camino de alto rendimiento. El wrapper \texttt{src/native.js} detecta si el binding está disponible y, si no, cae al motor JS. Esa decisión hace que el sistema sea portable durante desarrollo y de bajo riesgo durante despliegue.

\begin{lstlisting}[language=C++,caption={Esquema conceptual del motor en el addon nativo.}]
class ChlEngine {
public:
    uint32_t add(std::string_view text, Metadata meta);
    std::vector<Result> search(std::string_view query, size_t k);
    void reward(uint32_t query_id, uint32_t entry_id, float value);
    void compact();
};
\end{lstlisting}

La implementación prioriza:
\begin{itemize}
  \item \texttt{popcount} nativo;
  \item datos contiguos;
  \item menor asignación en el camino caliente;
  \item índices compactos;
  \item determinismo por semilla.
\end{itemize}

\section{Benchmarks reales}

Los benchmarks están definidos en \texttt{scripts/evaluate.js}. El modo \texttt{large} usa 12\,000 documentos y 6\,000 consultas. El modo \texttt{huge} usa 50\,000 documentos y 20\,000 consultas, pero sigue corriendo sobre el preset \texttt{large} del runtime. Los números siguientes son los últimos resultados reportados por el script.

\subsection{Benchmark large}

\begin{table}[h]
\centering
\small
\begin{tabular}{lrr}
\toprule
\textbf{Métrica} & \textbf{chl-native} & \textbf{chl-js} \\
\midrule
docs & 12\,000 & 12\,000 \\
queries & 6\,000 & 6\,000 \\
learnedConcepts & 16 & 16 \\
learnedPhrases & 18 & 18 \\
loadMs & 4\,451.15 & 3\,571.75 \\
recall@1 & 0.93917 & 0.93917 \\
recall@3 & 0.93983 & 0.93983 \\
mrr & 0.93950 & 0.93950 \\
medianLatencyUs & 452.21 & 453.90 \\
p95LatencyUs & 1\,255.92 & 1\,329.83 \\
collisionBuckets & 1\,672 & 1\,672 \\
maxBucketSize & 11 & 11 \\
\bottomrule
\end{tabular}
\caption{Resultado del benchmark \texttt{bench:large}.}
\end{table}

\begin{table}[h]
\centering
\small
\begin{tabular}{lr}
\toprule
\textbf{Categoría} & \textbf{recall@1} \\
\midrule
large-positive & 0.93438 \\
large-paraphrase & 0.93286 \\
large-negation & 0.98351 \\
\bottomrule
\end{tabular}
\caption{Desglose por categoría del benchmark large.}
\end{table}

\subsection{Benchmark huge}

\begin{table}[h]
\centering
\small
\begin{tabular}{lrr}
\toprule
\textbf{Métrica} & \textbf{chl-native} & \textbf{chl-js} \\
\midrule
docs & 50\,000 & 50\,000 \\
queries & 20\,000 & 20\,000 \\
learnedConcepts & 60 & 60 \\
learnedPhrases & 66 & 66 \\
loadMs & 123\,221.24 & 108\,347.83 \\
recall@1 & 0.80925 & 0.80905 \\
recall@3 & 0.80960 & 0.80960 \\
mrr & 0.80943 & 0.80933 \\
medianLatencyUs & 292.75 & 316.77 \\
p95LatencyUs & 863.54 & 948.79 \\
collisionBuckets & 2\,789 & 2\,789 \\
maxBucketSize & 9 & 9 \\
\bottomrule
\end{tabular}
\caption{Resultado del benchmark \texttt{bench:huge}.}
\end{table}

\begin{table}[h]
\centering
\small
\begin{tabular}{lrp{0.55\linewidth}}
\toprule
\textbf{Categoría} & \textbf{recall@1} & \textbf{Lectura} \\
\midrule
large-positive & 0.83913 & Buen comportamiento en consultas afirmativas directas. \\
large-paraphrase & 0.76454 & La paráfrasis sigue funcionando, pero ya muestra degradación al escalar. \\
large-negation & 0.86859 & La negación explícita se mantiene robusta. \\
\bottomrule
\end{tabular}
\caption{Desglose por categoría del benchmark huge.}
\end{table}

\subsection{Lectura de ingeniería}

Los números de los benchmarks dejan cuatro conclusiones realistas:
\begin{enumerate}
  \item La consulta es rápida incluso con decenas de miles de entradas.
  \item La carga inicial es el cuello de botella principal en 50k.
  \item La negación está bien modelada como señal explícita.
  \item La paráfrasis es la zona más sensible al crecimiento del corpus.
\end{enumerate}

\section{Comparación con alternativas de mercado}

\CHL{} no pretende reemplazar a Weaviate, Pinecone o Redis como motor universal de retrieval. La comparación correcta es de caso de uso.

\begin{table}[p]
\centering
\small
\begin{tabular}{p{0.17\linewidth}p{0.19\linewidth}p{0.19\linewidth}p{0.19\linewidth}p{0.19\linewidth}}
\toprule
\textbf{Criterio} & \textbf{CHL} & \textbf{Weaviate} & \textbf{Pinecone} & \textbf{Redis} \\
\midrule
Memoria de agente & Muy fuerte; conversacional y \MCP{}-native & Search generalista & Retrieval de producción & Infraestructura de datos \\
\MCP{} nativo & Sí & No & No & No \\
Explicabilidad & Alta: state, journal, backup, buckets y lexicón & Media & Media-baja & Media \\
Paráfrasis & Fuerte con lexicón aprendido & Buena con hybrid search & Buena con dense+sparse & Buena con vector + filtros \\
Negación & Señal explícita & Depende del ranking & Depende del esquema & Depende del índice \\
Local-first & Sí & Posible & Normalmente gestionado & Local o gestionado \\
Persistencia & Journal, backup y sidecars & Base de datos & Servicio & Infraestructura \\
Escala grande & Memoria local; no foco principal & Muy fuerte & Muy fuerte & Muy fuerte \\
Filtros/schema & Básico a medio & Muy fuerte & Fuerte & Muy fuerte \\
Coste operativo & Bajo & Medio/alto & Medio/alto & Medio \\
Mejor uso & Agente personal/producto & Búsqueda híbrida & RAG de producción & Search + cache + datos \\
\bottomrule
\end{tabular}
\caption{Comparación resumida con alternativas de mercado.}
\end{table}

La clave es simple: si el problema es memoria conversacional local e inspeccionable, \CHL{} tiene una ventaja clara. Si el problema es retrieval documental de gran escala, las plataformas generalistas siguen siendo superiores.

\section{Limitaciones reales}

\begin{itemize}
  \item \textbf{No hay contexto infinito:} la memoria crece, pero el coste de carga y compactación existe.
  \item \textbf{La paráfrasis cae al escalar:} el benchmark de 50k muestra degradación frente a 12k.
  \item \textbf{El modelo no reemplaza comprensión profunda:} la codificación de características sigue siendo decisiva.
  \item \textbf{La similitud filtra información:} \LSH{} no es privacidad.
\end{itemize}

\section{Conclusión}

El resultado final de \CHL{} es una memoria de agente local-first, \MCP{}-native, persistente, aprendible y explicable. El sistema ya no está en fase de idea: existe como implementación en C++ y JavaScript, con HTTP, \MCP{}, backup binario, journal, sidecars del lexicón, tests y benchmarks.

La contribución técnica más importante es que el sistema separa tres cosas que suelen mezclarse:
\begin{itemize}
  \item el direccionamiento semántico rápido por firmas binarias;
  \item la composición y limpieza por hipervectores;
  \item el vocabulario aprendido que persiste entre sesiones.
\end{itemize}

Los resultados experimentales son coherentes con esa lectura. En 12k, la paráfrasis supera 0.93 de recall@1; en 50k baja a 0.76, pero la consulta sigue siendo rápida y la negación se mantiene robusta. Eso sugiere que \CHL{} ya es útil como memoria de agente, pero que el siguiente trabajo debe enfocarse en compactación, escalado del lexicón y preservación de paráfrasis sin aumentar ruido.

\section*{Apéndice: interfaz mínima para uso real}

\begin{itemize}
  \item HTTP: \texttt{/lexicon}, \texttt{/lexicon.tsv}, \texttt{/backup.bin}, \texttt{/state}.
  \item \MCP{} tools: \texttt{chl\_remember}, \texttt{chl\_recall}, \texttt{chl\_infer}, \texttt{chl\_lexicon}, \texttt{chl\_backup\_binary}, \texttt{chl\_restore\_binary}.
  \item Recursos \MCP{}: \texttt{chl://memory}, \texttt{chl://lexicon}, \texttt{chl://backup.bin}.
\end{itemize}

\begin{thebibliography}{99}

\bibitem{indyk1998approximate}
P. Indyk y R. Motwani.
\newblock Approximate nearest neighbors: towards removing the curse of dimensionality.
\newblock \emph{Proceedings of the 30th Annual ACM Symposium on Theory of Computing}, 1998.

\bibitem{charikar2002similarity}
M. S. Charikar.
\newblock Similarity estimation techniques from rounding algorithms.
\newblock \emph{Proceedings of the 34th Annual ACM Symposium on Theory of Computing}, 2002.

\bibitem{broder1997resemblance}
A. Z. Broder.
\newblock On the resemblance and containment of documents.
\newblock \emph{Compression and Complexity of Sequences}, 1997.

\bibitem{kanerva1988sdm}
P. Kanerva.
\newblock \emph{Sparse Distributed Memory}.
\newblock MIT Press, 1988.

\bibitem{kanerva2009hdc}
P. Kanerva.
\newblock Hyperdimensional computing: An introduction to computing in distributed representation with high-dimensional random vectors.
\newblock \emph{Cognitive Computation}, 1(2):139--159, 2009.

\bibitem{plate1995hrr}
T. A. Plate.
\newblock Holographic reduced representations.
\newblock \emph{IEEE Transactions on Neural Networks}, 6(3):623--641, 1995.

\bibitem{gayler2003vsa}
R. W. Gayler.
\newblock Vector symbolic architectures answer Jackendoff's challenges for cognitive neuroscience.
\newblock \emph{Joint International Conference on Cognitive Science}, 2003.

\bibitem{mcpResources}
Model Context Protocol.
\newblock Resources.
\newblock Documentación oficial. \url{https://modelcontextprotocol.io/docs/concepts/resources}.

\bibitem{mcpTools}
Model Context Protocol.
\newblock Tools.
\newblock Documentación oficial. \url{https://modelcontextprotocol.io/specification/2025-06-18/server/tools}.

\bibitem{weaviateHybrid}
Weaviate.
\newblock Hybrid search.
\newblock Documentación oficial. \url{https://docs.weaviate.io/weaviate/concepts/search/hybrid-search}.

\bibitem{pineconeHybrid}
Pinecone.
\newblock Hybrid search.
\newblock Documentación oficial. \url{https://docs.pinecone.io/guides/search/hybrid-search}.

\bibitem{redisVector}
Redis.
\newblock Vector search concepts.
\newblock Documentación oficial. \url{https://redis.io/docs/latest/develop/ai/search-and-query/vectors/}.

\end{thebibliography}

\end{document}
